\documentclass[a4paper,12pt]{article}
\usepackage[T2A]{fontenc}
\usepackage[utf8]{inputenc}
\usepackage[russian]{babel}
\usepackage{amsmath,amssymb,amsfonts}
\usepackage{graphicx}
\usepackage{float}
\usepackage{caption}
\usepackage{subcaption}
\usepackage{geometry}
\geometry{left=3cm, right=1.5cm, top=2cm, bottom=2cm}
\usepackage{setspace}
\onehalfspacing % Полуторный интервал
\usepackage{indentfirst} % Красная строка
\setlength{\parindent}{1.25cm} % Отступ первой строки
\usepackage{enumitem}
\usepackage{booktabs} % Для красивых таблиц
\usepackage{hyperref}
\hypersetup{
    colorlinks=true,
    linkcolor=blue,
    filecolor=magenta,
    citecolor=red,     
    urlcolor=cyan,
    pdftitle={Курсовой проект},
    pdfpagemode=FullScreen,
}


\usepackage{amsthm} % Для теорем, лемм, определений
\usepackage{thmtools} % Для дополнительного оформления
\usepackage{xcolor} % Для цветов

% Определяем стили для окружений
\declaretheoremstyle[
    spaceabove=6pt, spacebelow=6pt,
    headfont=\normalfont\bfseries,
    notefont=\mdseries, notebraces={(}{)},
    bodyfont=\normalfont,
    postheadspace=1em,
    headpunct={}
]{mystyle}

% Создаем окружения с нумерацией внутри разделов
\declaretheorem[style=mystyle, name=Определение, numberwithin=section]{definition}
\declaretheorem[style=mystyle, name=Теорема, sibling=definition]{theorem}
\declaretheorem[style=mystyle, name=Лемма, sibling=definition]{lemma}
\declaretheorem[style=mystyle, name=Следствие, sibling=definition]{corollary}
\declaretheorem[style=mystyle, name=Утверждение, sibling=definition]{proposition}
\declaretheorem[style=mystyle, name=Пример, numbered=no]{example}



\begin{document}

% Титульный лист
\begin{titlepage}
    \centering
    {\large\bfseries Национальный исследовательский университет \\ «Высшая школа экономики»\par}
    \vspace{1cm}
    {\large Факультет экономических наук, Факультет компьютерных наук\par}
    {\large Образовательные программы «Экономика и анализ данных», «Прикладная математика и информатика»\par}
    \vspace{2cm}
    {\Large\bfseries КУРСОВОЙ ПРОЕКТ\par}
    \vspace{1cm}
    {\Large\bfseries Геометрические методы на многообразии матриц фиксированного ранга для заполнения матриц \par}
    \vspace{2cm}
    \begin{flushright}
        \large
        \textbf{Выполнили:}\\
        Студент 2 курса, группы БЭАД245\\
        Аляутдинов Карим Маратович,\\
        Студент 2 курса, группы БПМИИ245\\
        Исмаилов Магомед Айдынович.\\
        \vspace{0.5cm}
        \textbf{Научный руководитель:}\\
        профессор, д.ф.-м.н.\\
        Широков Дмитрий Сергеевич\\
    \end{flushright}
    \vfill
    {\large Москва, 2026\par}
\end{titlepage}

% Содержание
\tableofcontents
\newpage

% Основной текст
\section*{Введение}
\addcontentsline{toc}{section}{Введение}
\phantomsection

\textbf{Актуальность темы.} Здесь объясните, почему выбранная тема важна в современном экономическом и социальном контексте. Укажите на существующие проблемы или пробелы в знаниях, которые данная работа может частично заполнить. Актуальность должна быть подкреплена ссылками на современные тенденции, данные или литературу.

\textbf{Цель исследования.} Четко сформулируйте конечную цель работы (например: «Целью данной работы является оценка влияния факторов X, Y и Z на экономический показатель A с использованием методов анализа данных на примере региона B»).

\textbf{Задачи исследования.} Перечислите конкретные шаги для достижения цели.

\textbf{Объект и предмет исследования.}
\begin{itemize}
    \item \textbf{Объект:} Широкая область (например, рынок труда РФ, малые и средние предприятия, потребительское поведение).
    \item \textbf{Предмет:} Конкретные связи, свойства или процессы внутри объекта, которые будут изучаться (например, зависимость между уровнем образования и заработной платой).
\end{itemize}

\textbf{Методы исследования.} Перечислите используемые методы: анализ литературы, методы описательной статистики, методы визуализации данных, регрессионный анализ (укажите конкретные модели, например, метод наименьших квадратов, логистическая регрессия), методы проверки гипотез и другие методы.

\textbf{Структура работы.} Кратко опишите, из каких разделов состоит работа.

Объем введения -- не менее 2 страниц.


\newpage

\section{Обзор литературы}
\label{sec:literature}




Проанализируйте предыдущие исследования по вашей теме. Сравните их цели, данные, методы и основные выводы. Не забудьте корректно указать ссылки на источники. Помимо текста, можно также сделать обобщающую таблицу (в качестве примера, см. Таблицу \ref{tab:literature}).
\begin{table}[H]
    \centering
    \caption{Сводная таблица исследований по теме}
    \label{tab:literature}
    \begin{tabular}{p{0.15\linewidth}p{0.25\linewidth}p{0.2\linewidth}p{0.3\linewidth}}
        \toprule
        \textbf{Автор(ы), год}        & \textbf{Цель/Гипотеза} & \textbf{Данные и методы} & \textbf{Ключевые результаты}     \\
        \midrule
        Smith (2020) \cite{smith2020} & Оценить влияние A на B & Панельные данные, МНК    & Обнаружено положительное влияние \\
        \bottomrule
    \end{tabular}
\end{table}
Обзор литературы должен выглядеть не как перечисление работ с их кратким содержанием, а как критический анализ преимуществ, ограничений и пробелов работ предшественников и их сравнение. Опишите выбор методов/моделей/алгоритмов для Вашей работы.


\newpage

\section{Основные определения и обозначения}
\label{sec:definitions}

Введите основные опеределения, математические обозначения и сокращения, используемые в работе. Например:

\subsection*{Математические обозначения}

\begin{itemize}
    \item $\mathbb{R}$ --- множество действительных чисел.
    \item $\mathbb{R}^n$ --- $n$-мерное евклидово пространство.
    \item $\mathbf{X} = [x_{ij}]$ --- матрица, где $x_{ij}$ --- элемент, находящийся на пересечении $i$-й строки и $j$-го столбца.
    \item \ldots
\end{itemize}

\subsection*{Ключевые экономико-статистические понятия}

\begin{itemize}
    \item \textbf{Наблюдение (observation, data point)} --- запись в наборе данных, соответствующая одной единице исследования в определённый момент или за определённый период времени.
    \item \textbf{Признак (feature), переменная (variable)} --- измеряемая характеристика наблюдения.
    \item ...
\end{itemize}

\subsection*{Список сокращений}

\begin{itemize}
    \item \textbf{ВВП (GDP)} --- валовой внутренний продукт.
    \item \textbf{ВРП (GRP)} --- валовой региональный продукт.
    \item \textbf{МНК (OLS)} --- метод наименьших квадратов (Ordinary Least Squares).
    \item ...
\end{itemize}

\subsection*{Обозначения, специфичные для данной работы}

\begin{itemize}
    \item Для согласованности в данной работе приняты следующие обозначения основных переменных:
          \begin{itemize}
              \item $GDP\_growth_{it}$ --- годовой темп прироста валового внутреннего продукта региона $i$ в году $t$, \%.
              \item ...
          \end{itemize}
    \item ...
\end{itemize}


\newpage

\section{Теоретическая часть}
\label{sec:theory}

В этой части стоит разместить доказательства теорем, описание полученных / использованных в работе методов, моделей, алгоритмов. Текст должен быть написан строго, аккуратно и последовательно, как в математическом учебнике.

\subsection{Основные теоретические концепции}
Опишите ключевые теории, модели или концепции, лежащие в основе вашего исследования. Используйте математические формулы, если это уместно. Например,
\begin{eqnarray}\label{def:reg}
    Y_i = \beta_0 + \beta_1 X_{1i} + \beta_2 X_{2i} + \varepsilon_i,
\end{eqnarray}
где $Y_i$ — зависимая переменная, $X_{1i}, X_{2i}$ — значения признаков, $\beta_0, \beta_1, \beta_2$ — коэффициенты, $\varepsilon_i$ — случайная ошибка. На формулу (\ref{def:reg}) можно ссылаться.

Для тех, кто пока не знаком с \LaTeX и Overleaf, рекомендуем воспользоваться документацией Overleaf с примерами (см. \url{https://www.overleaf.com/learn/latex/Learn_LaTeX_in_30_minutes} и другие статьи их сайта).


\subsection{...}

Подпараграфы можно разбивать на подподпараграфы, если необходимо:
\subsubsection{...}

...

\subsubsection{...}

Ниже приведены примеры некоторых окружений, которые могут быть использованы в теоретической части работы (мы добавили их вручную в верхней части tex файла). Можно добавлять и другие окружения.

Напоминаем отличия между основными окружениями. \textit{Теорема (Theorem)} — главный, значительный результат, который требует доказательства. Это центральные утверждения работы.
\textit{Лемма (Lemma)} — вспомогательное утверждение, которое используется для доказательства теоремы. Часто это "технический"  результат.
\textit{Следствие (Corollary)} — результат, который непосредственно вытекает из теоремы, почти очевидное следствие.
\textit{Утверждение (Proposition)} — нечто среднее между леммой и теоремой, важное, но не центральное утверждение.



\begin{lemma}[Ортогональность собственных векторов симметричной матрицы]\label{lem:eigenvectors_orthogonal}
    Пусть $A \in \mathbb{R}^{n \times n}$ — симметричная матрица ($A = A^\top$). Тогда:
    \begin{enumerate}
        \item Все собственные значения $\lambda$ матрицы $A$ вещественны.
        \item Собственные векторы, соответствующие различным собственным значениям, ортогональны.
    \end{enumerate}
\end{lemma}

\begin{proof}
    Докажем первое утверждение ... Докажем второе утверждение ...
\end{proof}

\begin{theorem}[Спектральная теорема для симметричных матриц]\label{thm:spectral}
    Для любой симметричной матрицы $A \in \mathbb{R}^{n \times n}$ существует ортогональная матрица $Q$ (такая что $Q^\top Q = I$) и диагональная матрица $\Lambda$, для которых:
    \begin{eqnarray}
        A = Q \Lambda Q^\top.
    \end{eqnarray}
    При этом столбцы $Q$ являются собственными векторами $A$, а диагональные элементы $\Lambda$ — соответствующими собственными значениями.
\end{theorem}
\begin{proof}
    ... По Лемме \ref{lem:eigenvectors_orthogonal}, получаем ...
\end{proof}


\begin{corollary}[След и определитель через собственные значения]\label{cor:trace_det}
    Для симметричной матрицы $A$:
    \begin{enumerate}
        \item $\operatorname{tr}(A) = \sum_{i=1}^n \lambda_i$ (след равен сумме собственных значений)
        \item $\det(A) = \prod_{i=1}^n \lambda_i$ (определитель равен произведению собственных значений)
    \end{enumerate}
\end{corollary}

\begin{proof}
    Из спектрального разложения $A = Q\Lambda Q^\top$ имеем:
    \begin{enumerate}
        \item $\operatorname{tr}(A) = \operatorname{tr}(Q\Lambda Q^\top) = \operatorname{tr}(\Lambda Q^\top Q) = \operatorname{tr}(\Lambda) = \sum_{i=1}^n \lambda_i$
        \item $\det(A) = \det(Q\Lambda Q^\top) = \det(Q)\det(\Lambda)\det(Q^\top) = \det(\Lambda) = \prod_{i=1}^n \lambda_i$, так как $\det(Q) = \det(Q^\top) = \pm 1$.
    \end{enumerate}
\end{proof}


\begin{example}
    В методе главных компонент (PCA) для центрированной матрицы данных $X \in \mathbb{R}^{m \times n}$ (m наблюдений, n признаков) ковариационная матрица $C = \frac{1}{m-1}X^\top X$ является симметричной и положительно полуопределенной.
    По спектральной теореме существует разложение:
    \begin{eqnarray}
        C = Q\Lambda Q^\top,
    \end{eqnarray}
    где столбцы $Q$ — главные компоненты (направления максимальной дисперсии), а $\lambda_i$ — дисперсии вдоль этих направлений. По следствию \ref{cor:trace_det}, полная дисперсия данных равна:
    \begin{eqnarray}
        \operatorname{tr}(C) = \sum_{i=1}^n \lambda_i,
    \end{eqnarray}
    что позволяет вычислять долю объясненной дисперсии каждой главной компонентой.
\end{example}




\newpage

\section{Практическая часть}
\label{sec:empirical}

В этой части стоит разместить описание программной реализации / описание экспериментов, анализ полученных результатов. Код Вашей работы должен быть выложен в созданный Вами репозиторий на \href{https://github.com/}{github}, на его фрагменты можно будет ссылаться в тексте.


\subsection{Описание данных и источников}
\begin{itemize}
    \item \textbf{Источники данных:} ...
    \item \textbf{Период:} ...
    \item \textbf{Единицы наблюдения:} ...
    \item \textbf{Переменные:} Если актуально, представьте список переменных с их описанием, единицами измерения и ожидаемым влиянием (см. Таблицу \ref{tab:variables}).
\end{itemize}

\begin{table}[H]
    \centering
    \caption{Описание переменных}
    \label{tab:variables}
    \begin{tabular}{p{0.2\linewidth}p{0.4\linewidth}p{0.15\linewidth}}
        \toprule
        \textbf{Переменная}  & \textbf{Описание}                    & \textbf{Единица измерения} \\
        \midrule
        \textit{gdp\_growth} & Темп роста ВРП                       & \%                         \\
        \textit{invest}      & Объем инвестиций в основной капитал  & млн руб.                   \\
        \textit{edu\_level}  & Доля населения с высшим образованием & \%                         \\
        \bottomrule
    \end{tabular}
\end{table}

\subsection{Предобработка и описательный анализ}
Опишите шаги по очистке данных: работа с пропущенными значениями, выбросами, кодированием категориальных переменных и др. Приведите основные описательные статистики (например, среднее, медиана, стандартное отклонение, минимум, максимум) в таблице. Добавьте графики. На графиках должны быть крупно подписаны оси и значения на осях, графики должны быть понятны читателю. На каждый график нужно сослаться в тексте и пояснить его (что на нем изображено, какие из этого следуют выводы и т.д.).
\begin{figure}[H]
    \centering
    % \includegraphics[width=0.7\linewidth]{plot_reg_5.png} % Замените на свой график
    \caption{Значения функции потерь (MSE) на тестовой выборке (ось OY) в зависимости от размера обучающей выборки (ось OX) на моделях EMLP-O(5), EMLP-SO(5), MLP, MLP+Aug, CGENN, GLGENN (наша модель).}
    \label{fig:dist}
\end{figure}

\subsection{Модель}
Сформулируйте спецификацию вашей модели (например, сама модель, ее составляющие, выбранные гиперпараметры и т.д.). Обоснуйте выбор метода. Если это модель машинного обучения, опишите детали ее обучения (например, оптимизатор, его гиперпараметры и т.д.). Обоснуйте выбор. Опишите метрики, по которым планируется сравнивать результаты, и обоснуйте их выбор. Если уместно, можно построить схему работы Вашей модели (например, схему с основными блоками, составляющими модель, и входными/выходными данными; схему взаимодействия пользователя с Вашей моделью и т.д.).


\subsection{Результаты}

\textbf{Результаты.} Представьте результаты в виде таблиц (например, как Таблица \ref{tab:regression}), графиков (lineplot, barplot, histogram и другие), рисунков.

\begin{table}[H]
    \centering
    \caption{Результаты оценки регрессионной модели}
    \label{tab:regression}
    \begin{tabular}{lccc}
        \toprule
        \textbf{Переменная}  & \textbf{Модель 1} & \textbf{Модель 2} & \textbf{Модель 3 (FE)} \\
        \midrule
        invest               & 0.45***           & 0.38***           & 0.29**                 \\
                             & (0.12)            & (0.11)            & (0.13)                 \\
        edu\_level           &                   & 0.21**            & 0.15*                  \\
                             &                   & (0.09)            & (0.08)                 \\
        Константа            & 1.05***           & 0.82***           & 0.91***                \\
                             & (0.30)            & (0.28)            & (0.31)                 \\
        \midrule
        N                    & 85                & 85                & 85                     \\
        R\textsuperscript{2} & 0.65              & 0.72              & 0.78                   \\
        \bottomrule
        \multicolumn{4}{l}{\footnotesize{*** p<0.01, ** p<0.05, * p<0.1; стандартные ошибки в скобках.}}
    \end{tabular}
\end{table}

\textbf{Интерпретация результатов.} Подробно объясните значения результатов и их значимость. Сравните результаты. Сделайте выводы.


\newpage

\section*{Заключение}
\addcontentsline{toc}{section}{Заключение}
\phantomsection

\textbf{Личный вклад.} Выделите Ваш личный вклад (например, <<в разделе XXX описана предложенная автором модификация алгоритма>>).

\textbf{Выводы.} Кратко и последовательно изложите основные выводы по каждой из задач, поставленных во введении. Избегайте повторения результатов, обобщайте их.

\textbf{Ограничения исследования.} Честно укажите на ограничения вашей работы (например: ограниченный период данных и т.д.).



\textbf{Приложения результатов}. Расскажите о предполагаемых применениях полученных результатов.


\textbf{Направления для дальнейших исследований.} Предложите, как можно развить вашу работу в будущем (использование других методов машинного обучения, данных большего периода и т.д.).

Объем заключения -- не менее 2 страниц.

\newpage

\begin{thebibliography}{99}
    \addcontentsline{toc}{section}{Список литературы}
    \phantomsection

    \bibitem{wooldridge}
    Wooldridge, J. M. (2016). Introductory Econometrics: A Modern Approach. Cengage Learning.

    \bibitem{rosstat}
    Федеральная служба государственной статистики (Росстат). \url{https://rosstat.gov.ru}

    \bibitem{smith2020}
    Smith, J. (2020). The Impact of Education on Economic Growth. Journal of Economic Perspectives, 34(4), 123-145.

    \bibitem{rlms}
    Российский мониторинг экономического положения и здоровья населения НИУ ВШЭ (RLMS-HSE). \url{https://rlms.hse.ru}

    % Добавьте свои источники здесь

\end{thebibliography}

\newpage

% Приложения (если нужны)
\section*{Приложение А. Дополнительные таблицы и графики}
\addcontentsline{toc}{section}{Приложение А. Дополнительные таблицы и графики}
\phantomsection
\setcounter{section}{1}
\setcounter{table}{0}
\renewcommand{\thetable}{А.\arabic{table}}
\setcounter{figure}{0}
\renewcommand{\thefigure}{А.\arabic{figure}}

Здесь можно разместить, например:
\begin{itemize}
    \item Дополнительные таблицы с полными результатами экспериментов.
    \item Дополнительные графики.
    \item Важные фрагменты кода.
    \item Опросные листы (если использовались).
\end{itemize}

\end{document}